\documentclass[11pt]{article}
\usepackage[utf8]{vietnam}
\usepackage[
    paperwidth=8.5in,
    paperheight=11in,
    top=4.4in,
    bottom=0.85in,
    left=0.85in,
    right=0.85in
]{geometry}

\usepackage{amsmath,amssymb}
\usepackage{graphicx}
\usepackage{xcolor}
\usepackage{fancyhdr}
\usepackage{tikz}

% Bảng màu đặc trưng của giáo trình James Stewart Calculus
\definecolor{stewartcyan}{RGB}{0, 130, 185}
\definecolor{stewartdark}{RGB}{35, 35, 35}
\definecolor{creditgray}{RGB}{115, 115, 115}

% Cấu hình đầu trang và chân trang
\pagestyle{fancy}
\fancyhf{}
\renewcommand{\headrulewidth}{0pt}
\renewcommand{\footrulewidth}{0pt}
\fancyfoot[R]{\fontsize{11}{13}\selectfont\sffamily\bfseries\color{stewartdark}279}

\begin{document}
\thispagestyle{fancy}

% 4 vạch ngang trang trí mép trái và hình ảnh đầu chương tràn lề trên-phải
\begin{tikzpicture}[remember picture, overlay]
    % 4 vạch ngang màu xanh cyan ở mép trái trang giấy
    \foreach \i in {0, 1, 2, 3} {
        \draw[stewartcyan, line width=1.1pt] 
            ([xshift=0pt, yshift={-0.85in - \i*4.5pt}] current page.north west) -- 
            ([xshift=0.85in, yshift={-0.85in - \i*4.5pt}] current page.north west);
    }
    % Ảnh đầu chương tràn lề đỉnh (top) và lề phải (right)
    \node[anchor=north west, inner sep=0pt] at ([xshift=0.85in] current page.north west) {
        \includegraphics[width=7.65in, height=4.25in]{images/p0314_fig1.png}
    };
\end{tikzpicture}%

% Chú thích ảnh và trích dẫn của Leonard Euler
\noindent
{\fontsize{8.5}{12}\selectfont\sffamily\bfseries\color{stewartdark}
Nhà toán học vĩ đại Leonard Euler từng nhận định: ``... trong vũ trụ không có điều gì diễn ra mà ở đó không xuất hiện một quy tắc nào đó về cực đại hoặc cực tiểu.'' Trong Bài tập 4.7.53, bạn sẽ sử dụng giải tích để chỉ ra rằng loài ong xây dựng các ô trong tổ của chúng theo hình dạng tối thiểu hóa diện tích bề mặt.
\par}

\vspace{0.35em}

% Ghi công hình ảnh
\noindent
{\fontsize{7}{9}\selectfont\sffamily\color{creditgray}
Kostiantyn Kravchenko / Shutterstock.com
\par}

\vspace{2.2em}

% Tiêu đề chương: Số 4 lớn và tên chương dóng hàng đáy (baseline)
\noindent
\begin{minipage}[b]{1.4in}
    \raggedleft
    {\fontsize{74}{74}\selectfont\bfseries\sffamily\color{stewartcyan}4}
\end{minipage}%
\hspace{16pt}%
% Đường phân cách dọc màu xanh cyan kéo dài xuống dưới đoạn văn

\begin{tikzpicture}[overlay]
    \draw[stewartcyan, line width=1.4pt] (0pt, 54pt) -- (0pt, -225pt);
\end{tikzpicture}%
\hspace{16pt}%
\begin{minipage}[b]{\dimexpr\textwidth - 1.4in - 32pt\relax}
    {\fontsize{26}{30}\selectfont\bfseries\sffamily\color{stewartdark}Các ứng dụng của đạo hàm\par}%
\end{minipage}

\vspace{1.8em}

% Đoạn văn mở đầu chương dóng thẳng với tên chương
\noindent
\hspace*{\dimexpr 1.4in + 32pt\relax}%
\begin{minipage}[t]{\dimexpr\textwidth - 1.4in - 32pt\relax}
    {\fontsize{10.5}{15}\selectfont\color{stewartdark}
    \textbf{CHÚNG TA ĐÃ TÌM HIỂU MỘT SỐ} ứng dụng của đạo hàm, nhưng giờ đây khi đã nắm vững các quy tắc tính đạo hàm, chúng ta có vị thế thuận lợi hơn để đi sâu nghiên cứu các ứng dụng của đạo hàm một cách toàn diện hơn. Tại đây, chúng ta sẽ học xem đạo hàm cho chúng ta biết điều gì về hình dạng đồ thị của một hàm số và đặc biệt là cách chúng giúp ta xác định các giá trị cực đại và cực tiểu của hàm số. Nhiều bài toán thực tế đòi hỏi chúng ta phải tối thiểu hóa chi phí, tối đa hóa diện tích, hoặc bằng cách nào đó tìm ra kết quả tối ưu nhất cho một tình huống. Cụ thể, chúng ta sẽ có thể khảo sát hình dạng tối ưu của một vỏ lon và giải thích vị trí xuất hiện của cầu vồng trên bầu trời.
    \par}
\end{minipage}

\end{document}